Nodes that share most of their neighbours tend to share a community. On this premise,
three nodes, one adjacent to the other two, whose three pairs all reach a commonality
threshold form an elementary community, and communities are the connected components of
the resulting 3-uniform hypergraph under shared pairs. The construction, HyperCommon, is
deterministic, local and nested in its threshold, and every edge it admits lies in
exactly one community. A randomly rewired edge rarely joins two nodes with many
neighbours in common, so the threshold rejects it while the triples inside each group
hold the group together. Compared with sixteen other methods, every parameter chosen
without labels, HyperCommon keeps $\Omega\ge0.6$ to a higher rewired fraction than every
other method in thirteen of twenty-eight ring-lattice settings as the rings are rewired
from order to randomness; its lead holds under a paired bootstrap in nine, all with
rings of unequal size or preferential rewiring (on one, up to $0.68$ against $0.63$ for
the next method). It is the most accurate in sixteen of the settings and level in four,
and in ten when every method gets its best parameter on each graph. We derive the exact
window of thresholds that recovers planted overlapping ring lattices, and show that once
ring degrees differ enough no single threshold recovers both their dense--sparse and
their sparse--sparse overlaps. One threshold fits rings of 600 and of 30 nodes alike: on
such a gradient without noise HyperCommon, like clique percolation, link communities and
ego-splitting, reaches $\Omega\ge0.97$ at its best parameter, Leiden at most $0.32$ at
its best resolution. It trails on all eight families of overlapping LFR graphs, and on
the rings once a tenth (with labels) or three tenths (without) of their edges are
rewired inside the rings: it needs communities whose members share many neighbours. On
real networks, including eleven kinds built to have lattice-like groups, it is among the
three most accurate methods on four kinds but rarely the most accurate, because a single
threshold cannot cut the strongest links between groups and keep the weakest inside
them.
